% =============================================================================
% SECTION 8: CONNECTION TO CMET
% =============================================================================
\section{Connection to CMET}\label{sec:cmet}

This manuscript is the second of a planned trilogy. The first paper (CMET, currently in revision) establishes the \emph{manifold-extension primitive}: for a single foundation MLIP $M$, the set of configurations on which $M$'s prediction error exceeds a threshold forms a smooth submanifold of the regular configuration space under suitable regularity conditions. CMET proves that this error manifold has finite reach, bounded intrinsic dimension, and stable homology under sampling.

The present paper promotes the CMET primitive from a single-model statement to a \emph{class-level} universality theorem. The logical relationship is:
\begin{equation}\label{eq:cmet-universality}
    \underbrace{\text{CMET}}_{\text{single-model manifold}} \xrightarrow{\text{Cross-Model Vandermonde}} \underbrace{\text{Universality Theorem}}_{\text{class-uniform geometry}}.
\end{equation}

Specifically, CMET provides the following primitives that the Universality Theorem extends:
\begin{itemize}
    \item \textbf{Manifold existence.} CMET proves that $\cM(M)$ is a smooth submanifold for any $M$ satisfying (F3). The Universality Theorem bounds $\dim \cM(M)$ uniformly over $M \in \cF$ via the Vandermonde lemma.
    \item \textbf{Reach bound.} CMET gives a single-model reach bound depending on $M$'s second-derivative bound. The Universality Theorem replaces this with $\inf_M \reach(\cM(M)) \geq r(\cF)$, a class-uniform constant.
    \item \textbf{Sample complexity.} CMET's single-model sample complexity is $n \geq C(M) \cdot d(M) \log(N/d(M))$. The Universality Theorem replaces $C(M)$ and $d(M)$ with class constants $C(\cF)$ and $d(\cF)$.
\end{itemize}

The Cross-Model Vandermonde Lemma (Lemma~\ref{lem:vandermonde}) is the mathematical engine that enables this promotion. Where CMET takes the intrinsic dimension $d(M)$ and Fisher decay rate $\rho(M)$ as model-dependent quantities, the Vandermonde lemma proves that both are uniformly bounded across $\cF$. The proof assembles the single-model Vandermonde bound of Waterfall \emph{et al.}\cite{waterfall2006sloppy} with the class-uniform smoothness of (F3) and the class-uniform expressivity of (F1)--(F2) to close the loop.

The third paper in the trilogy will treat the \emph{refusal-mode} behavior of foundation MLIPs in full generality, building on the two-mode inference framework of clause (v). Where the present paper characterizes the refusal region $\cX_N^{\text{refuse}}(M)$ geometrically (via Federer reach), the third paper will develop the operational criteria for a model to refuse prediction---uncertainty quantification thresholds, out-of-distribution detection, and human-in-the-loop protocols for high-stakes applications such as materials discovery and drug design.

\begin{remark}
The meta-theorem framing---CMET as the primitive, the Universality Theorem as the promotion---is not merely organizational. It reflects a substantive logical dependency: the class-uniform constants $d(\cF)$, $\rho(\cF)$, and $r(\cF)$ cannot be proved without the single-model regularity established by CMET. The Universality Theorem is, in this sense, a \emph{conditional} theorem: it holds for any class $\cF$ whose members individually satisfy the CMET regularity conditions, provided the class-level uniformity conditions (F1)--(F3) are also satisfied.
\end{remark}
